% =============================================================
% 2026 年高教社杯 C 题《微网与外部电网电力调控策略》
% 题目原文中全部 4 个表格的 LaTeX 代码汇总
%
% 使用说明：
%   1. 表格含中文，请使用 ctex 文档类或在导言区加载 \usepackage{ctex}，
%      并使用 XeLaTeX 编译；
%   2. 以下 4 个表格相互独立，可直接整段复制到论文中使用；
%   3. 表格中的空白单元格为原题预留的填写位置；
%   4. 除 ctex（中文支持）外无需额外宏包。
% =============================================================


% -------------------------------------------------------------
% 表 1  微网在指定时间段的购电量及全天的购电量和购电费（单位：kWh、元）
% -------------------------------------------------------------
\begin{table}[htbp]
  \centering
  \caption{微网在指定时间段的购电量及全天的购电量和购电费}
  \label{tab:1}
  \begin{tabular}{|c|c|c|c|c|c|}
    \hline
    时间段      & 购电量 & 时间段      & 购电量 & 时间段      & 购电量 \\ \hline
    10:00-10:10 &        & 12:00-12:10 &        & 14:00-14:10 &        \\ \hline
    16:00-16:10 &        & 18:00-18:10 &        & 20:00-20:10 &        \\ \hline
    \multicolumn{2}{|c|}{全天购电量} & &
    \multicolumn{2}{c|}{全天购电费} &        \\ \hline
  \end{tabular}
\end{table}


% -------------------------------------------------------------
% 表 2  储能设备在指定时间段的充放电量及 0:00 和 24:00 的储电量（单位：kWh）
% -------------------------------------------------------------
\begin{table}[htbp]
  \centering
  \caption{储能设备在指定时间段的充放电量及 0:00 和 24:00 的储电量}
  \label{tab:2}
  \begin{tabular}{|c|c|c|c|c|c|}
    \hline
    时间段     & 充电量 & 放电量 & 时间段      & 充电量 & 放电量 \\ \hline
    0:00-4:00  &        &        & 4:00-8:00   &        &        \\ \hline
    8:00-12:00 &        &        & 12:00-16:00 &        &        \\ \hline
    16:00-20:00&        &        & 20:00-24:00 &        &        \\ \hline
    \multicolumn{2}{|c|}{0:00 储电量}  & &
    \multicolumn{2}{c|}{24:00 储电量} &        \\ \hline
  \end{tabular}
\end{table}


% -------------------------------------------------------------
% 表 3  微网在指定日期的紧急购电量（单位：kWh）
% 列数较多，局部使用 \small 并缩小列间距以保证不超出版心
% -------------------------------------------------------------
\begin{table}[htbp]
  \centering
  \caption{微网在指定日期的紧急购电量}
  \label{tab:3}
  \small
  \setlength{\tabcolsep}{4pt}
  \begin{tabular}{|c|c|c|c|c|c|c|c|}
    \hline
    \multicolumn{2}{|c|}{2025.3.20}  & \multicolumn{2}{c|}{2025.6.21} &
    \multicolumn{2}{c|}{2025.9.23}  & \multicolumn{2}{c|}{2025.12.21} \\ \hline
    时间段 & 购电量 & 时间段 & 购电量 & 时间段 & 购电量 & 时间段 & 购电量 \\ \hline
            &        &        &        &        &        &        &        \\ \hline
            &        &        &        &        &        &        &        \\ \hline
            &        &        &        &        &        &        &        \\ \hline
  \end{tabular}
\end{table}


% -------------------------------------------------------------
% 表 4  紧急购电填写示例（单位：kWh）
% 原题中日期格内的横线是贯通的，故按原文留空处理；
% 若希望日期“2025/3/1”在三行中垂直居中，可将第一列改为
% \multirow{3}{*}{2025/3/1}（需 \usepackage{multirow}），
% 并把行间的 \hline 换成 \cline{2-3}。
% -------------------------------------------------------------
\begin{table}[htbp]
  \centering
  \caption{紧急购电填写示例}
  \label{tab:4}
  \begin{tabular}{|c|c|c|}
    \hline
    日期      & 紧急购电时间段 & 紧急购电量 \\ \hline
    2025/3/1  & 13:00-13:30    & 100        \\ \hline
              & 14:40-15:50    & 400        \\ \hline
              & 16:10-16:30    & 50         \\ \hline
  \end{tabular}
\end{table}
